\documentclass[11pt]{article}
\usepackage[a4paper,margin=1in]{geometry}
\usepackage[T1]{fontenc}
\usepackage[utf8]{inputenc}
\usepackage{lmodern}
\usepackage{amsmath,amssymb,amsthm}
\usepackage{enumitem}
\usepackage{microtype}
\usepackage{hyperref}
\newcommand{\webersection}[2]{\section*{\S #1. #2}\addcontentsline{toc}{section}{\S #1. #2}}
\newtheorem*{satz}{Theorem}
\title{Heinrich Weber\\\textit{Lehrbuch der Algebra}, Volume II\\Current English Translation Draft}
\author{}
\date{}
\begin{document}
\maketitle

\section*{Reader note}
This current working draft combines the available English translation pieces for Heinrich Weber's \emph{Lehrbuch der Algebra}, Volume II: the preface, \S\S 1--8 of the general group theory opening, and the current continuations on abelian groups, bases, invariants, group characters, reciprocal groups, Lagrange resolvents, linear substitutions, transposed substitutions, projective actions, binary substitutions, absolute invariants, form problems, Galois groups, metacyclic groups, abstract groups, alternating groups, orthogonal groups, and fractional linear groups.
The translation is intended as a continuation base for checking and refinement, not as a proofed edition.

\section*{Preface to the Second Volume}
In accordance with the intention announced in the preface to the first volume, I am now able to place the second volume of my \emph{Textbook of Algebra} before the public. The plan laid down there has been carried out in its essential points. In the applications I have tried to select problems that had already acquired an independent interest in other fields, such as geometry or the theory of functions, and that at the same time bring the principal points of algebraic theory before the reader in as many different ways as possible.

The application of the theory of algebraic numbers has been carried through as far as the theory of cyclotomic numbers. If life and strength endure, I hope in a continuation of my work to present the further applications to the domain of elliptic functions, which are contained only in part in my book \emph{Elliptic Functions and Algebraic Numbers}.

During the preparation and printing of the second volume as well, I have had beside me the help and counsel of the friends whom I already named in the preface to the first edition. But the first volume has by now also won for itself many new friends who have advanced my work by suggestions and advice. To them all I express my thanks here, and add the request that they may continue to preserve their interest in the work.

\begin{flushright}
Strasbourg, July 1896.\\
The Author.
\end{flushright}

\section*{Part I. General Group Theory}
\section*{\S 1. Definition of Groups}
In the first volume, in connection with permutations, we became acquainted with the concept of a group and made important algebraic applications of it. Our next task must therefore be to formulate more generally this concept, so extraordinarily important in all of modern mathematics, and to become familiar with the laws that govern it. We place the following definition at the head of the discussion.

A system $P$ of objects (elements) of any kind whatsoever becomes a \emph{group} when the following assumptions are satisfied.

\begin{enumerate}
\item A rule is given according to which, from a first and a second element of the system, a completely determined third element of the same system is derived.
\end{enumerate}

Symbolically, if $a$ is the first, $b$ the second, and $c$ the third element, one writes
\[
ab=c,\qquad c=ab,
\]
and says that $c$ is \emph{composed} from $a$ and $b$, while $a$ and $b$ are the \emph{components} of $c$.

For this composition one does not in general assume the commutative law; that is, $ab$ may differ from $ba$. On the other hand, one does assume

\begin{enumerate}
\setcounter{enumi}{1}
\item the \emph{associative law},
\end{enumerate}
that is, whenever $a$, $b$, $c$ are any three elements of $P$, one has
\[
(ab)c=a(bc).
\]

From this, by complete induction, it follows that one always arrives at the same result when, in any finite sequence of elements of $P$, $a$, $b$, $c$, $d$, \dots, one first composes two neighboring elements, then again two neighboring elements, and so on until the whole sequence has been reduced to a single element, denoted by $abcd\dots$. Thus, for example,
\begin{align*}
abcd &= (ab)cd = [(ab)c]d = (ab)(cd) \\
     &= a(bc)d = [a(bc)]d = a[(bc)d] \\
     &= ab(cd) = (ab)(cd) = a[b(cd)].
\end{align*}

\bigskip

\clearpage

\section*{§2. Intersection of two groups}
Cosets do not themselves have the defining properties of a group. For, in order that the composition of two elements of $Q_1$ should again yield an element of $Q_1$, one would need, for example,
\[
 a_1 b\, a_2 b = a_3 b,
\]
and hence
\[
 a_1 b a_2 = a_3,
 \qquad
 b = a_1^{-1} a_3 a_2^{-1}.
\]
Thus $b$ would, contrary to the assumption, have to belong to $Q$. The designation \emph{coset} is therefore to be understood only in an inexact or analogical sense.

Two divisors $Q,Q'$ of a group $P$ always have the element $1$ in common. They may, however, have further elements in common, and these common elements form a group, which we denote by $R$. For if $a$ and $b$ belong both to $Q$ and to $Q'$, then by the group property of $Q$ and $Q'$ the same is true of the product $ab$; therefore $R$ is a group. This common group is called the greatest common divisor of $Q$ and $Q'$, or, with a geometric turn of phrase, the \emph{intersection} of $Q$ and $Q'$.

\section*{§3. Normal subgroups}
If $Q$ is a divisor of $P$ that is identical with all of its conjugate divisors, then $Q$ is called a \emph{normal subgroup} of $P$.

The group consisting only of the single element $1$ is a normal subgroup of every group. We obtain another normal subgroup by taking the greatest common divisor $R$ of all divisors conjugate to some divisor $Q$ of $P$.

Indeed, it has already been proved that $R$, as the intersection of several divisors of $P$, is itself a group. Now let
\[
 Q,\;Q',\;Q'',\dots
\]
be the divisors conjugate to $Q$, and let $b$ be any element of $P$. Then the system of groups
\[
 b^{-1}Qb,\; b^{-1}Q'b,\; b^{-1}Q''b,\dots
\]
is not different from the original system. If $R$ is the intersection of the first system, then $b^{-1}Rb$ is the intersection of the second system, and consequently $R=b^{-1}Rb$; that is, $R$ is a normal subgroup of $P$.

If $N$ is a normal subgroup of a group $P$ and $b$ is any element of $P$, then
\[
 b^{-1} N b = N,
 \qquad\text{or equivalently}\qquad
 N b = b N.
\]
If $P$ has order $n$, $N$ has order $\nu$, and index $\mu$, so that $n=\mu\nu$, then we may choose the $\mu$ elements $1,b_1,\dots,b_{\mu-1}$ in such a way that the decomposition of $P$ into cosets is
\[
 P = N + N b_1 + N b_2 + \cdots + N b_{\mu-1}
   = N + b_1 N + b_2 N + \cdots + b_{\mu-1} N.
\]

\section*{§4. Composition of subsets}
Let $A$ and $B$ be any two collections of elements of a group $P$ (whether groups or not). By the symbolic product $AB$ we understand the set of all elements obtained by composing each element $a$ of $A$ with each element $b$ of $B$.

The two equations
\[
 AA = A,
 \qquad
 AB = BA,
\]
where $A$ is fixed and $B$ is an arbitrary subset of $P$, express that $A$ is a normal subgroup of $P$. For from the first relation $A$ is itself a divisor of $P$, and from the second it follows for every element $b\in P$ that
\[
 b^{-1} A b = A,
\]
that is, $A$ is normal in $P$.

If $A$ and $B$ are two divisors of $P$, then one of the two equalities
\[
 AB = BA,
 \qquad\text{or equivalently}\qquad
 BA = AB,
\]
is the necessary and sufficient condition for the product set $AB$ to be a group.

Because of the defining property of normal subgroups, $N$ commutes with every $b$, and therefore
\[
 b_\nu N = N b_\nu,
\]
so that
\[
 N b_\nu \cdot N b_\mu = N N \cdot b_\nu b_\mu.
\]
Since $b_\nu b_\mu$ lies in $P$, the set $N b_\nu b_\mu$ coincides with one of the cosets occurring in the decomposition above; and because $N N = N$, it follows that if we write $N b_\nu b_\mu = M$, then
\[
 N M = M.
\]
Thus property (1) in the definition of a group has been verified for the composition of cosets. That the associative law also holds was already noted. It remains only to verify the analogue of cancellation.

\section*{§5. Multiple isomorphism}
These quotient groups are related to one another in the same way as isomorphic groups in the ordinary sense, except that to each element $A$ there may correspond not just one element $a$, but several. This gives occasion to extend the notion of isomorphism by the following definition.

A group $P$ with elements $a,b,c,\dots$ is said to be \emph{multiply isomorphic} to a group $Q$ with elements $A,B,C,\dots$ if the two groups can be put into correspondence in such a way that to each of the elements $A,B,C,\dots$ there correspond one or more of the elements $a,b,c,\dots$, while every element of $Q$ corresponds to one and only one element of $P$, and such that whenever $a$ corresponds to $A$ and $b$ corresponds to $B$, then $ab$ corresponds to $AB$.

In this way one proves first that to each of the elements of the quotient group there corresponds an entire coset of the original group. Following H"older, one writes such a quotient group in the form $P/N$. The order of the complementary group is equal to the index of the group $N$.

We mention in conclusion a theorem about normal subgroups. If $N$ is a divisor of $P$, and $N'$ is a divisor of $N$ which is at the same time a normal subgroup of $P$, then $N'$ is also a normal subgroup of $N$. This is immediate, because for every element $b\in P$ one has
\[
 b^{-1}N'b = N',
\]
and every element of $N$ is also an element of $P$.

The converse, however, is not valid in general: if $N$ is a normal subgroup of $P$ and $N'$ is a normal subgroup of $N$, it need not follow that $N'$ is a normal subgroup of $P$.

\section*{§6. Composition series and the theorem of C. Jordan}
A group $P$ is called \emph{simple} if it has no normal subgroup other than itself and the identity; it is called \emph{composite} if other normal subgroups are present.

A normal subgroup which is not contained properly in any larger proper normal subgroup of $P$ is called a \emph{maximal normal subgroup} of $P$. If $P_1$ is a maximal normal subgroup of $P$, $P_2$ is a maximal normal subgroup of $P_1$, and so on, then one obtains a chain
\[
 P,\; P_1,\; P_2,\; \dots,\; P_\nu = 1,
\]
in which each subgroup is a maximal normal subgroup of the preceding one. Such a chain is called a \emph{composition series}; the corresponding indices are its indices of composition.

The important theorem expressing the uniqueness of the composition factors was established by C.~Jordan.

\medskip
\noindent\textbf{Theorem (Jordan).} \emph{However a composition series of a group $P$ may be chosen, the sequence of its indices is always the same up to order.}
\medskip

The proof rests on the quotient groups introduced above. Suppose $Q$ and $Q_1$ are normal subgroups of $P$, and that in addition $Q_1$ is a divisor (hence, by the previous section, a normal subgroup) of $Q$. Then the quotient group $Q/Q_1$ is a normal subgroup of $P/Q_1$, and the index of $Q/Q_1$ in $P/Q_1$ is equal to the index of $Q$ in $P$.

This reduction is the crucial device in the proof that the composition factors are independent of the chosen composition series, apart from their order.

\clearpage

\section*{\S 7. Further results on composition series}

We derive first several important consequences of Jordan's theorem.

Let $P$ be a group and $Q$ a normal subgroup of $P$, and suppose that $P$ and $Q$ have composition factors
\[
A_1,A_2,\dots,A_\lambda
\qquad\text{and}\qquad
C_1,C_2,\dots,C_\mu.
\]
Then the quotient group $P/Q$ has composition factors
\[
A_1,A_2,\dots,A_\lambda,\; C_\mu^{-1},\dots,C_2^{-1},C_1^{-1}.
\]
For every composition series of $P/Q$ can be completed to one of $P$.

\medskip
\noindent\textbf{Corollary 1.1.}
If $P,Q,R$ are three groups, each of which is a normal subgroup of the preceding one, and if the index of $R$ in $P$ is a power of a prime number, then the index of $R$ in $Q$ and the index of $Q$ in $P$ are powers of the same prime.

\section*{\S 8. Metacyclic groups}

A group $P$ is called \emph{metacyclic} if it possesses a composition series all of whose factors are cyclic.

This definition includes all cyclic groups, and also all groups whose order is a power of a prime number. For every group of prime-power order is, by Sylow's theorem, solvable, hence metacyclic.

Metacyclic groups play an important role in the theory of algebraic equations, since the equations whose groups are metacyclic are solvable by radicals.

\section*{Appendix: foundations of group theory}

The preceding paragraphs contain the foundations of general group theory as they are used in modern algebra. We have treated the definition of a group, the theory of divisors and cosets, normal subgroups and their properties, many-valued isomorphism, and composition series together with Jordan's theorem.

These concepts are of fundamental importance for the further development of algebra. Upon them is based the theory of abelian groups, which is presented in the following paragraphs of the second section.

\clearpage

\section*{Second Section. Abelian groups}

\section*{\S 9. Representation of Abelian groups by a basis}

We have already called those groups in which the commutative law holds for composition \emph{commutative} or \emph{Abelian} groups. For these groups, which are of particular importance in applications, the laws are much simpler than in general groups, as we now see.

Let \(S\) be an Abelian group of order \(n\). We shall show that the elements of \(S\) can be represented in the form
\[
\theta=A_1^{\alpha_1}A_2^{\alpha_2}\cdots A_\nu^{\alpha_\nu},
\tag{1}
\]
where \(A_1,\dots,A_\nu\) are suitable elements of the group and the exponents \(\alpha_1,\dots,\alpha_\nu\) run through complete residue systems modulo positive integers \(a_1,\dots,a_\nu\), the orders of \(A_1,\dots,A_\nu\). Such a system \(A_1,\dots,A_\nu\) is called a \emph{basis} of the Abelian group.

One sees immediately that every element of the form \((1)\) actually occurs: for example, one obtains \(A_i\) by taking \(\alpha_i=1\) and \(\alpha_k=0\) for \(k\neq i\).

It remains to prove that every element occurs equally often. Suppose
\[
1=A_1^{h_1}A_2^{h_2}\cdots A_n^{h_n}
\tag{2}
\]
is a representation of the identity. Then the value of \(\theta\) in \((1)\) is unchanged if one replaces \(\alpha_1,\dots,\alpha_n\) by \(\alpha_1+h_1,\dots,\alpha_n+h_n\). Hence each element \(\theta\) is represented at least as many times by \((1)\) as the identity is represented by \((2)\).

Conversely, if two exponent systems \(\alpha_1,\dots,\alpha_n\) and \(\alpha_1',\dots,\alpha_n'\) yield the same element \(\theta\), then
\[
1=A_1^{\alpha_1-\alpha_1'}A_2^{\alpha_2-\alpha_2'}\cdots A_n^{\alpha_n-\alpha_n'}.
\tag{3}
\]
Thus the difference of the two exponent systems gives a representation of the identity. Therefore the number of representations of \(\theta\) can be no larger than the number of representations of \(1\). If we denote that number by \(h\), then the total number of possible exponent choices in \((1)\) is \(a_1a_2\cdots a_n\), while \(n\) is the number of distinct group elements. Hence
\[
nh=a_1a_2\cdots a_n.
\tag{4}
\]

From formula \((4)\) one obtains:

\medskip
\noindent\textbf{Proposition 2.}
If \(r\) is a prime dividing the order \(n\) of the group \(S\), then \(S\) contains an element of order \(r\).

\medskip
Indeed, one of the orders \(a_1,\dots,a_n\) must be divisible by \(r\). If, say, \(a_1=rk\), then \(A_1^k\) has order \(r\).

\medskip
\noindent\textbf{Proposition 3.}
If \(A,B,\dots\) are arbitrary elements of \(S\) with orders \(a,b,\dots\), then the product \(AB\cdots\) again lies in \(S\), and its order divides the least common multiple \(m\) of \(a,b,\dots\).

\medskip
For
\[
(AB\cdots)^m=A^mB^m\cdots=1,
\]
so \(m\) is a multiple of the order of \(AB\cdots\).

\medskip
\noindent\textbf{Proposition 4.}
Suppose the order \(n\) of the group is factored as
\[
n=ab
\]
with \(a\) and \(b\) relatively prime. Then there exist exactly \(a\) elements \(A\) whose order divides \(a\), and exactly \(b\) elements \(B\) whose order divides \(b\), and every element of the group is represented uniquely in the form
\[
\theta=AB.
\tag{5}
\]

To see this, let \(\mathcal A\) be the set of elements whose order divides \(a\). By Proposition 3 this is a subgroup of \(S\). Similarly the set \(\mathcal B\) of elements whose order divides \(b\) is a subgroup.

Let \(a'\) be the order of \(\mathcal A\), and \(b'\) the order of \(\mathcal B\). Then \(a'\) is relatively prime to \(b\): if a prime \(r\) divided \(a'\), then by Proposition 2 the subgroup \(\mathcal A\) would contain an element of order \(r\), and since every element of \(\mathcal A\) has order dividing \(a\), that prime \(r\) would divide \(a\), not \(b\). Similarly \(b'\) is relatively prime to \(a\).

Now choose integers \(x,y\) such that
\[
ax+by=1.
\tag{6}
\]
For any element \(\theta\in S\) we have
\[
\theta=\theta^{ax}\theta^{by}.
\tag{7}
\]
Since \((\theta^{ax})^b=1\), the element \(\theta^{ax}\) lies in \(\mathcal A\), and similarly \(\theta^{by}\in\mathcal B\). Therefore every element has a representation
\[
\theta=AB.
\tag{8}
\]

This representation is unique. If \(AB=A'B'\), then by raising to the suitable powers \(by=1-ax\) and \(ax=1-by\) one obtains \(A=A'\) and hence \(B=B'\). The number of products \(AB\) is therefore \(a'b'\), and since every element of \(S\) occurs once, we have
\[
n=ab=a'b'.
\]
Because \(a\) is relatively prime to \(b'\) and \(b\) is relatively prime to \(a'\), it follows that
\[
a=a',\qquad b=b'.
\]
This proves Proposition 4 in full.

If the two subgroups \(\mathcal A,\mathcal B\) are each represented by bases, then formula \((8)\) shows that \(S\) itself is represented by a basis, namely by adjoining the basis elements of \(\mathcal A\) and \(\mathcal B\).

Repeating this decomposition whenever \(a\) and \(b\) can themselves be split into relatively prime factors, we reduce matters to the case where the order of the group is a power of a single prime:
\[
n=p^\mu.
\tag{9}
\]

Thus it remains to prove that every Abelian group of prime-power order can be represented by a basis. Let \(A\) be an element of maximal order \(a\). Then \(a\) is a power of \(p\), and the order of every element of \(S\) divides \(a\). Hence for every \(\theta\in S\),
\[
\theta^a=1.
\tag{10}
\]

The elements
\[
1,A,A^2,\dots,A^{a-1}
\tag{11}
\]
are all distinct and form a subgroup of \(S\). If they already exhaust the group, then \(S\) is represented by the one-element basis \(A\).

Otherwise, for every \(\theta\in S\) there exists a least positive integer \(b\) such that
\[
\theta^b=A^t
\tag{12}
\]
for some integer \(t\). This \(b\) divides \(a\), hence is a power of \(p\), and it also divides the order of \(\theta\). From \((12)\) one finds that \(t\) must be divisible by \(b\): writing \(a=qb+b'\) with \(0\leq b'<b\), one checks that if \(b'\neq 0\) then \(\theta^{b'}\) would already be a power of \(A\), contradicting the minimality of \(b\).

Writing \(a/b=k\) and setting
\[
B=\theta A^{-tk},
\tag{13}
\]
we obtain
\[
B^b=1,
\tag{14}
\]
and \(b\) is the least exponent with this property. Thus the elements
\[
A^{\alpha_1}B^{\alpha_2},
\qquad
\alpha_1=0,1,\dots,a-1,\quad
\alpha_2=0,1,\dots,b-1,
\tag{15}
\]
are all distinct and form a subgroup \(S'\) represented by the basis \(A,B\). If \(S'=S\), we are done; otherwise the construction continues inductively.

Proceeding by induction, one obtains a subgroup \(S_{\nu-1}\) with basis
\[
S_{\nu-1}=A_1^{\alpha_1}A_2^{\alpha_2}\cdots A_{\nu-1}^{\alpha_{\nu-1}},
\tag{16}
\]
such that:
\begin{enumerate}
\item the orders satisfy
\[
a_1\geq a_2\geq \cdots\geq a_{\nu-1},
\]
\item for every element \(\theta\in S\), the power \(\theta^{a_{\nu-1}}\) lies in \(S_{\nu-1}\), and the corresponding exponents are all divisible by \(a_{\nu-1}\).
\end{enumerate}

If \(S_{\nu-1}\neq S\), choose \(\theta\in S\) so that the least positive exponent \(a_\nu\) with \(\theta^{a_\nu}\in S_{\nu-1}\) is as large as possible. Then \(a_\nu\leq a_{\nu-1}\), and \(a_\nu\) is again a power of \(p\). Writing
\[
\theta^{a_\nu}=A_1^{h_1a_\nu}A_2^{h_2a_\nu}\cdots A_{\nu-1}^{h_{\nu-1}a_\nu},
\tag{17}
\]
define
\[
A_\nu=\theta A_1^{-h_1}A_2^{-h_2}\cdots A_{\nu-1}^{-h_{\nu-1}}.
\tag{18}
\]
Then
\[
A_\nu^{a_\nu}=1,
\tag{19}
\]
and \(a_\nu\) is the smallest such exponent. Consequently, the elements
\[
A_1^{\alpha_1}A_2^{\alpha_2}\cdots A_\nu^{\alpha_\nu},
\qquad
0\leq \alpha_i<a_i,
\tag{20}
\]
are all distinct and form a larger subgroup \(S_\nu\), again satisfying the inductive conditions.

Continuing this process, one eventually obtains the whole group \(S\) represented by a basis. This proves the theorem that every finite Abelian group can be represented by a basis.

It follows at once that every Abelian group is metacyclic. Indeed, if
\[
\theta=A_1^{\alpha_1}A_2^{\alpha_2}\cdots A_\nu^{\alpha_\nu},
\]
and if \(p\) is a prime dividing the order \(a_1\) of \(A_1\), then the elements
\[
\theta'=A_1^{p\alpha_1}A_2^{\alpha_2}\cdots A_\nu^{\alpha_\nu}
\]
form a subgroup of index \(p\), represented by the basis \(A_1^p,A_2,\dots,A_\nu\). Repeating this procedure shows that \(S\) admits a chain of subgroups of prime index, hence is metacyclic.

\section*{\S 10. The invariants of Abelian groups}

The proof just given for the existence of a basis also provides a method for finding one in which the orders of the basis elements are prime powers. Nevertheless, the same Abelian group can admit different bases.

For example, let \(A\) and \(B\) have relatively prime orders \(a\) and \(b\). Then they determine a two-element basis with group
\[
A^\alpha B^\beta,
\qquad
\alpha=0,1,\dots,a-1,\quad \beta=0,1,\dots,b-1.
\tag{1}
\]
But the same group can also be represented by the one-element basis \(AB\), namely in the form
\[
(AB)^s,
\qquad
s=0,1,\dots,ab-1.
\tag{2}
\]
Indeed, given arbitrary \(\alpha,\beta\), one can determine \(s\) modulo \(ab\) so that
\[
s\equiv \alpha \pmod a,\qquad s\equiv \beta \pmod b,
\]
and then \((2)\) reproduces \((1)\).

Even so, there is a precise sense in which the structure of the basis is completely determined by the group itself.

\medskip
\noindent\textbf{Theorem 0.2.}
Suppose
\[
A_1,A_2,\dots,A_\nu
\]
with orders
\[
a_1,a_2,\dots,a_\nu
\]
and
\[
B_1,B_2,\dots,B_\mu
\]
with orders
\[
b_1,b_2,\dots,b_\mu
\]
are two bases of the same Abelian group \(S\) of order \(n\). Let \(p\) be a prime dividing \(n\), and write
\[
a_i=p^{\alpha_i}\bar a_i,\qquad b_j=p^{\beta_j}\bar b_j,
\]
where \(\bar a_i,\bar b_j\) are not divisible by \(p\). Then the prime powers \(p^{\alpha_i}\) greater than \(1\) occur among the \(p^{\beta_j}\), and conversely.

\medskip
To prove this, order the basis elements so that
\[
a_1\leq a_2\leq \cdots\leq a_\nu,\qquad
b_1\leq b_2\leq \cdots\leq b_\mu.
\tag{3}
\]
Let \(m\) be the least common multiple of \(a_1,\dots,a_\nu\). Since the \(\bar a_i\) are not divisible by \(p\), the number \(m\) is also not divisible by \(p\).

For an arbitrary element
\[
\theta=A_1^{\alpha_1}A_2^{\alpha_2}\cdots A_\nu^{\alpha_\nu},
\tag{4}
\]
one then has
\[
\theta^{p^{\alpha_\nu}m}=1.
\]
Applying this to the basis elements \(B_1,\dots,B_\mu\), one sees that every \(b_j\) divides \(p^{\alpha_\nu}m\). Hence \(p^{\beta_\mu}\) divides \(p^{\alpha_\nu}\), and the least common multiple of the \(b_j\) is again \(m\). Reversing the rôles of the two bases, one obtains the opposite divisibility, and therefore
\[
p^{\alpha_\nu}=p^{\beta_\mu}.
\tag{5}
\]

Assume inductively that
\[
p^{\alpha_1}=p^{\beta_1},\quad
p^{\alpha_2}=p^{\beta_2},\quad
\dots,\quad
p^{\alpha_{s-1}}=p^{\beta_{s-1}}
\tag{6}
\]
have already been proved. Then for any element \(\theta\),
\[
\theta^{p^{\alpha_s}m}
=
A_1^{\alpha_1p^{\alpha_s}m}
A_2^{\alpha_2p^{\alpha_s}m}
\cdots
A_\nu^{\alpha_\nu p^{\alpha_s}m}.
\tag{7}
\]
Counting the number of distinct elements of this form yields
\[
q=
\frac{p^{\alpha_1}}{p^{\alpha_s}}
\cdot
\frac{p^{\alpha_2}}{p^{\alpha_s}}
\cdots
\frac{p^{\alpha_{s-1}}}{p^{\alpha_s}}.
\tag{9}
\]
If one now expresses the same elements in the \(B\)-basis and assumes
\[
p^{\alpha_s}>p^{\beta_s},
\tag{10}
\]
then the counting argument forces \(p^{\beta_s}\) to divide \(p^{\alpha_s}\), which contradicts \((10)\) unless
\[
p^{\alpha_s}=p^{\beta_s}.
\tag{14}
\]
This proves Theorem 0.2.

The prime powers occurring in the orders of a basis are therefore independent of the choice of basis. We call them the \emph{invariants} of the Abelian group. Their product is equal to the order \(n\) of the group.

By \S 9 we can choose a basis whose element orders are prime powers. Hence those orders are exactly the invariants and are completely determined by the group.

That the invariants completely determine the group as well follows from the next theorem.

\medskip
\noindent\textbf{Theorem 0.3.}
Two Abelian groups with the same invariants are isomorphic, and isomorphic Abelian groups have the same invariants.

\medskip
Indeed, if
\[
\theta=A_1^{\alpha_1}A_2^{\alpha_2}\cdots A_\nu^{\alpha_\nu}
\]
are the elements of one group \(S\), and
\[
\theta'=B_1^{\alpha_1}B_2^{\alpha_2}\cdots B_\nu^{\alpha_\nu}
\]
those of another group \(S'\), with the same exponent ranges modulo the same basis orders \(a_1,\dots,a_\nu\), then we obtain an isomorphism by letting \(\theta\) correspond to \(\theta'\) whenever the exponent systems are identical.

Conversely, if two groups are isomorphic, then the images of a basis of one group form a basis of the other, and therefore the invariants must agree in both.

\clearpage

\section*{\S 11. Group characters}

A principal problem in the theory of Abelian groups is to determine all divisors of an Abelian group. The solution of this problem is substantially facilitated by introducing the concept of \emph{group characters}, which is also important in many other respects.

Let \(S\) be an Abelian group of order \(n\), whose elements we denote by \(A\), and let
\[
A_1,A_2,\dots,A_\nu
\tag{1}
\]
be the elements of a basis of \(S\), of orders
\[
a_1,a_2,\dots,a_\nu.
\]
Then every element \(A\) is represented, and represented uniquely, in the form
\[
A = A_1^{\alpha_1}A_2^{\alpha_2}\cdots A_\nu^{\alpha_\nu},
\tag{2}
\]
where
\[
0\le \alpha_1<a_1,\qquad 0\le \alpha_2<a_2,\qquad \dots,\qquad 0\le \alpha_\nu<a_\nu.
\tag{3}
\]
The exponents \(\alpha_1,\alpha_2,\dots,\alpha_\nu\), or any integers congruent to them modulo \(a_1,a_2,\dots,a_\nu\), are called the \emph{indices} of the element \(A\). Then one has:

\medskip
\noindent\textbf{1.} The indices of a product \(AA'\) are obtained by adding the corresponding indices of the two factors.

\medskip
If we assign to each element \(A\) some numerical value, we may call this assignment a function of \(A\). Such a function \(\chi(A)\) will be called a \emph{group character} if for every two elements \(A,A'\in S\) it satisfies
\[
\chi(AA')=\chi(A)\chi(A').
\tag{4}
\]
Hence it also satisfies the more general multiplicative law
\[
\chi(A'A''A'''\cdots)=\chi(A')\chi(A'')\chi(A''')\cdots.
\]

Two such functions \(\chi\) and \(\chi_1\) are regarded as distinct if there exists at least one element \(A\) for which \(\chi(A)\neq \chi_1(A)\). We now determine these characters more closely.

Setting \(A'=1\) in \((4)\), one obtains
\[
\chi(A)=\chi(A)\chi(1).
\tag{5}
\]
Excluding once and for all the degenerate case in which all \(\chi(A)=0\), one gets:
\[
\chi(1)=1.
\tag{6}
\]

\medskip
\noindent\textbf{2.} Every character takes the value \(1\) on the identity element.

\medskip
Setting \(A'=A\) in \((4)\), we obtain
\[
\chi(A^2)=[\chi(A)]^2,
\]
and by complete induction, for every exponent \(h\),
\[
\chi(A^h)=[\chi(A)]^h.
\tag{7}
\]
If \(a\) denotes the order of the element \(A\), then with \(h=a\) and using \((6)\) one gets
\[
[\chi(A)]^a=1.
\tag{8}
\]

\medskip
\noindent\textbf{3.} The values of a character \(\chi(A)\) are roots of unity whose orders divide the order of \(A\).

\medskip
We can now determine all characters easily. Set
\[
\chi(A_1)=\omega_1,\qquad \chi(A_2)=\omega_2,\qquad \dots,\qquad \chi(A_\nu)=\omega_\nu.
\tag{9}
\]
Then
\[
\omega_1^{a_1}=1,\qquad \omega_2^{a_2}=1,\qquad \dots,\qquad \omega_\nu^{a_\nu}=1,
\tag{10}
\]
and from \((2)\) and \((4)\),
\[
\chi(A)=\omega_1^{\alpha_1}\omega_2^{\alpha_2}\cdots \omega_\nu^{\alpha_\nu}.
\tag{11}
\]
Conversely, whenever \(\omega_1,\omega_2,\dots,\omega_\nu\) is any solution of \((10)\), formula \((11)\) defines a function of \(A\) that satisfies \((4)\), and hence is a character of the group.

Each of the equations \((10)\) has \(a_1,a_2,\dots,a_\nu\) roots, and if all root choices are combined with one another, one obtains
\[
a_1a_2\cdots a_\nu=n
\]
possible systems. These lead to \(n\) distinct characters \(\chi(A)\). Indeed, if another choice
\[
\omega'_1,\omega'_2,\dots,\omega'_\nu
\]
yields the same function for all \(A\), then taking \(\alpha_1=1\), \(\alpha_2=\cdots=\alpha_\nu=0\) shows \(\omega_1=\omega'_1\), and similarly for the others.

Thus:

\medskip
\noindent\textbf{4.} There exist exactly \(n\) distinct characters of an Abelian group of order \(n\), all given by formula \((11)\).

\medskip
Let \(\varepsilon_1,\varepsilon_2,\dots,\varepsilon_\nu\) be primitive roots of the equations \((10)\). Then we may write
\[
\omega_1=\varepsilon_1^{\beta_1},\qquad \omega_2=\varepsilon_2^{\beta_2},\qquad \dots,\qquad \omega_\nu=\varepsilon_\nu^{\beta_\nu},
\tag{12}
\]
where \(\beta_1,\dots,\beta_\nu\) are taken from complete residue systems modulo \(a_1,\dots,a_\nu\). Then all \(n\) characters are given by
\[
\chi(A)=\varepsilon_1^{\alpha_1\beta_1}\varepsilon_2^{\alpha_2\beta_2}\cdots\varepsilon_\nu^{\alpha_\nu\beta_\nu}.
\tag{13}
\]
Among these occurs the character obtained by setting all \(\beta_i=0\); it has value \(1\) on every element \(A\), and we call it the \emph{unit character}.

The \(n\) characters of \(S\) can themselves be combined into an Abelian group, indeed into one isomorphic to \(S\).

For by \((13)\), each character is determined by a system of integers
\[
\beta_1,\beta_2,\dots,\beta_\nu
\]
taken modulo \(a_1,\dots,a_\nu\), and this same system is the index system of a definite element \(B\in S\), namely
\[
B=A_1^{\beta_1}A_2^{\beta_2}\cdots A_\nu^{\beta_\nu}.
\tag{14}
\]
Thus every element \(B\in S\) corresponds to a definite character, which we denote by \(\chi_B(A)\), or simply \(\chi_B\) when the variable \(A\) is omitted.

If \(B'\) is a second element of \(S\), with indices \(\beta'_1,\beta'_2,\dots,\beta'_\nu\), then its corresponding character \(\chi_{B'}\) is defined likewise. If now \(\chi_B\chi_{B'}\) denotes the character given on each \(A\) by
\[
\chi_B\chi_{B'}(A)
=
\varepsilon_1^{\alpha_1(\beta_1+\beta'_1)}
\varepsilon_2^{\alpha_2(\beta_2+\beta'_2)}
\cdots
\varepsilon_\nu^{\alpha_\nu(\beta_\nu+\beta'_\nu)},
\tag{15}
\]
then the characters are thereby combined into a group isomorphic to \(S\). The identity element of this character group is the unit character.

Accordingly, the character group can itself be represented by a basis, obtained by setting
\[
\chi_1(A)=\varepsilon_1^{\alpha_1},\qquad
\chi_2(A)=\varepsilon_2^{\alpha_2},\qquad
\dots,\qquad
\chi_\nu(A)=\varepsilon_\nu^{\alpha_\nu}.
\tag{16}
\]
Then every character has the form
\[
\chi_B = \chi_1^{\beta_1}\chi_2^{\beta_2}\cdots \chi_\nu^{\beta_\nu},
\tag{17}
\]
and if \(B,B'\in S\), then
\[
\chi_B\chi_{B'}=\chi_{BB'}.
\tag{18}
\]
If \(\chi_1,\chi_2\) are any two characters and \(\chi_1\chi_2\) their composite, then for every \(A\),
\[
\chi_1(A)\chi_2(A) = (\chi_1\chi_2)(A).
\tag{19}
\]

We may summarize this as:

\medskip
\noindent\textbf{5.} The characters of a group \(S\) can be combined into a group isomorphic to \(S\).

\medskip
One has also, immediately from definition \((13)\),
\[
\chi_B(A)=\chi_A(B).
\tag{20}
\]

Finally there is the following theorem.

\medskip
\noindent\textbf{6.} If \(A\) runs through all elements of the group \(S\), then for a fixed character \(\chi\),
\[
\sum_A \chi(A)=
\begin{cases}
n,& \text{if }\chi\text{ is the unit character},\\[4pt]
0,& \text{otherwise}.
\end{cases}
\tag{21}
\]
Likewise, if the element \(A\) is held fixed and \(\chi\) runs through all characters, then
\[
\sum_\chi \chi(A)=
\begin{cases}
n,& \text{if }A=1,\\[4pt]
0,& \text{otherwise}.
\end{cases}
\tag{22}
\]

When \(\chi\) is the unit character, or \(A\) is the identity, these formulas are evident, since each of the \(n\) summands is then equal to \(1\). If \(\chi\) is not the unit character, there exists \(B\in S\) such that \(\chi(B)\neq 1\). Writing
\[
\sum_A \chi(A)=\sigma,
\]
and multiplying by \(\chi(B)\), one obtains
\[
\sum_A \chi(AB)=\sigma\chi(B).
\]
But as \(A\) runs through the whole group, so does \(AB\), hence the left side is also \(\sigma\). Therefore
\[
\sigma[1-\chi(B)]=0,
\]
and since \(\chi(B)\neq 1\), we get \(\sigma=0\), proving \((21)\). Formula \((22)\) follows similarly, and also immediately from \((20)\).

\section*{\S 12. Divisors of an Abelian group. Reciprocal groups}

A divisor \(T\) of an Abelian group \(S\) is, as already remarked above, always a normal subgroup. Hence by \S 4 there exists a complementary group
\[
S/T,
\]
whose elements are the cosets of \(T\), namely
\[
T,\quad TA',\quad TA'',\quad \dots,
\tag{1}
\]
where \(A',A'',\dots\) are certain elements of \(S\). If \(t\) is the order of \(T\), then the number of such cosets, that is, the order of the group \(S/T\), is equal to the index of \(T\) in \(S\), namely
\[
\frac{n}{t}.
\]

This group \(S/T\) is itself Abelian, because
\[
TA'\cdot TA'' = TA'A'',
\qquad
TA''\cdot TA' = TA''A',
\]
and these are equal.

Let us now determine the characters of the group \(S/T\). Denote its elements by
\[
T_0,T_1,T_2,\dots,T_j,
\tag{2}
\]
and let one of its characters be written \(\xi(T_i)\).

From this function \(\xi(T)\) we can derive a function \(\xi(A)\) on the elements of \(S\) by setting
\[
\xi(A_i)=\xi(T_i)
\tag{3}
\]
whenever \(A_i\) is any element belonging to the coset \(T_i\). For the elements of the group \(T\) itself one then has \(\xi(A)=1\).

But this function \(\xi(A)\) is among the characters of \(S\). For if \(A_i\in T_i\) and \(A_k\in T_k\), then \(A_iA_k\in T_iT_k\), and hence
\[
\xi(A_i)\xi(A_k)=\xi(T_i)\xi(T_k)=\xi(T_iT_k)=\xi(A_iA_k).
\tag{4}
\]
This is precisely the defining property of a character of \(S\).

Conversely, suppose one of the characters \(\chi=\xi\) of \(S\) takes the same value on every element of \(T\). Since the identity belongs to \(T\), this common value can only be \(1\). If now \(A_i\) is any element of the coset \(T_i\), and \(A\) runs through the whole group \(T\), then \(AA_i\) runs through the coset \(T_i\). Since \(\xi(A)=1\), one has
\[
\xi(AA_i)=\xi(A_i).
\tag{5}
\]
That is, \(\xi(A)\) has one and the same value on all elements of a fixed coset \(T_i\), and may therefore be regarded as a function of \(T_i\), denoted \(\xi(T_i)\).

Since whenever \(A_i\in T_i\) and \(A_k\in T_k\), the product \(A_iA_k\in T_iT_k\), it follows that
\[
\xi(T_i)\xi(T_k)=\xi(T_iT_k).
\tag{6}
\]
So \(\xi(T_i)\) is indeed a character of the quotient group \(S/T\).

Therefore there are exactly \(j\) such characters \(\xi(A)\), namely those that take the value \(1\) on every element of \(T\). These \(j\) characters form, under composition of characters, a group; for if \(\xi_1(A)=1\) and \(\xi_2(A)=1\) on \(T\), then also \(\xi_1\xi_2(A)=1\). Since these functions \(\xi\) can equally well be viewed as the characters of the quotient group \(S/T\), it follows from \S 11, theorem 5, that the group of the \(\xi\) is isomorphic to \(S/T\).

Hence:

\medskip
\noindent\textbf{7.} If an Abelian group \(S\) has a divisor \(T\) of order \(t\) and index \(j\), then among the characters of \(S\) there are exactly \(j\), and no more, that take the value \(1\) on all elements of \(T\); and for every element of \(S\) not contained in \(T\), at least one of these characters takes a value different from \(1\). These characters form a group isomorphic to \(S/T\).

\medskip
Since every character \(\chi_B\) corresponds to a definite element \(B\in S\), if \(\chi_B\) runs through the group of characters belonging to \(T\), then \(B\) runs through a group of order \(j\), isomorphic to the group of the \(\chi_B\), and therefore also to \(S/T\). We denote this group by \(U\) and call it the \emph{group reciprocal to \(T\)}.

Again, one should note that in general the notion of reciprocal group depends on the chosen basis and the chosen roots of unity \(\varepsilon\), so it does not belong intrinsically to the group \(S\) and the subgroup \(T\) alone.

The reciprocal group of \(T\) is characterized by the following theorem.

\medskip
\noindent\textbf{8.} Let \(A\) run through the elements of a divisor \(T\subset S\). Then the elements \(B\in S\) satisfying
\[
\chi_B(A)=1
\tag{7}
\]
for every \(A\in T\) are precisely the elements of the group \(U\) reciprocal to \(T\), whose order is equal to the index of \(T\).

\medskip
By \S 11, \((20)\), the group \(T\) is in turn the reciprocal group to \(U\); the relation between the two groups is therefore mutual.

One also has the following theorem.

\medskip
\noindent\textbf{9.} If \(T'\) is a divisor of \(T\), then conversely the reciprocal group \(U\) of \(T\) is a divisor of the group \(U'\) reciprocal to \(T'\).

\medskip
For every element \(A'\in T'\) also lies in \(T\). Hence, if \(B\) runs through \(U\), then \(\chi_B(A')=1\), and therefore \(B\) must occur in the group \(U'\).

If one selects from the totality of characters \(\chi\) an arbitrary finite number
\[
\xi_1,\xi_2,\dots,\xi_k,
\]
whether or not they form a group, then all elements \(A\) satisfying the \(k\) conditions
\[
\xi_1(A)=1,\qquad
\xi_2(A)=1,\qquad
\dots,\qquad
\xi_k(A)=1
\tag{8}
\]
form a group, because from \(\xi_i(A)=1\) and \(\xi_i(A')=1\) it follows that \(\xi_i(AA')=1\).

If the characters \(\xi_1,\dots,\xi_k\) do not already form a group, then there follow from equations \((8)\) additional equations
\[
\xi_{k+1}(A)=1,\quad \dots
\]
so that \(\xi_1,\dots,\xi_k\) may be completed to a group.

\medskip
\noindent\textbf{10.} It follows from theorem 7 that in this way one can obtain all possible divisors of \(S\).

\medskip

\clearpage

\section*{\S 4. Cyclic equations}

\subsection*{163. Reduction of Abelian equations to cyclic equations}

We have already seen above that in a transitive Abelian group, besides the identity, no permutation occurs that leaves any symbol fixed. From this one obtains another theorem.

Assume that a permutation $\pi$ of such a group is decomposed into its cycles, and let the cycle with the fewest elements be an $r$-cycle. Then $\pi^r$ leaves the elements of this cycle fixed and therefore must be the identity permutation. But from $\pi^r = 1$ it follows that, if other cycles are present, they too must have $r$ elements. Thus:

\medskip
\noindent\textbf{Theorem.}
A permutation of a transitive Abelian group contains only cycles of the same length.

\medskip
If $m$ is the degree of the group, then the number $r$ of elements in one cycle must divide $m$, and if $m = rs$, then $s$ is the number of $r$-cycles of which $\pi$ consists.

Now let $P$ be the group of an Abelian equation $F(x)=0$, and choose any non-identity permutation $\pi$ in $P$. Decompose it into cycles
\[
\pi = \gamma_1\gamma_2\cdots\gamma_s,
\tag{1}
\]
where each cycle $\gamma$ refers to $r$ roots of the equation $F(x)=0$. Arrange these roots so that
\[
\gamma_1=(\alpha,\alpha_1,\alpha_2,\dots,\alpha_{r-1}),
\]
\[
\gamma_2=(\beta,\beta_1,\beta_2,\dots,\beta_{r-1}),
\]
\[
\dots,
\]
\[
\gamma_s=(\delta,\delta_1,\delta_2,\dots,\delta_{r-1}).
\tag{2}
\]

If $\pi_i$ is any permutation of $P$, then because of commutativity one has
\[
\pi^{-1}\pi_i\pi=\pi_i,
\qquad\text{that is,}\qquad
\pi_i\pi=\pi\pi_i.
\tag{3}
\]
By the theorem on the formation of the permutations $\pi^{-1}\pi_i\pi$ (\S 154, no.\ 6), it follows that $\gamma_1$ cannot be changed when the permutations $\pi$ are carried out. But since $\gamma_1$ is completely determined inside the cycles of $\pi$, apart from its ordering and its initial element, it follows that under any permutation $\pi_i$ in $P$ the elements belonging to one cycle $\gamma$ cannot be separated from one another; they can only be cyclically permuted among themselves, and the cycles can moreover be permuted among one another.

Hence, if $s>1$, the group is imprimitive.

A rational function of the arguments $\alpha,\alpha_1,\dots,\alpha_{r-1}$ that does not change its value when the $\alpha$ are subjected to the cyclic permutation $\gamma_1$ and its repetitions is called a \emph{cyclic function} of the $\alpha$. If we denote by $\omega,\omega_1,\dots,\omega_{s-1}$ the conjugate values of $\omega$, setting for instance
\[
\omega=\psi(\alpha,\alpha_1,\dots,\alpha_{r-1}),
\]
\[
\omega_1=\psi(\beta,\beta_1,\dots,\beta_{r-1}),
\qquad\dots,
\tag{4}
\]
then these quantities are, by \S 158, the roots of an irreducible equation of degree $s$,
\[
\Phi(\omega)=0,
\tag{5}
\]
whose group is obtained by determining the permutations induced on the quantities \((4)\) by the permutations coming from $P$.

But the permutation of the quantities \((4)\) induced by $\pi_i\pi_j$ is obtained by composing the permutations induced separately by $\pi_i$ and $\pi_j$; hence the group of permutations of \((4)\) is itself commutative.

Therefore $\Phi(\omega)=0$ is an Abelian equation of degree $s$. Its first factor $F(x,\omega)$ has the roots $\alpha,\alpha_1,\dots,\alpha_{r-1}$, and the group of the equation
\[
F(x,\omega)=0
\]
consists only of the permutations belonging to the period of the cyclic permutation $\gamma_1$.

We shall call an equation whose group consists of a single cycle and its repetitions a \emph{cyclic equation}. Thus cyclic equations are the simplest special case of Abelian equations.

We have therefore proved that the solution of every Abelian equation is reduced to the solution of an Abelian equation of lower degree together with the solution of a collection of cyclic equations.

One can apply the same theorem again to the auxiliary equation of degree $s$, and continue in this way until that auxiliary equation is reduced to degree one.

Thus we obtain the result:

\medskip
\noindent\textbf{Theorem.}
The solution of an Abelian equation can always be reduced to the solution of a sequence of cyclic equations whose degrees divide the degree of the given equation.

\medskip
It is not necessary to include irreducibility in the definition of cyclic equations. We may therefore formulate the definition in general as follows.

An equation of degree $m$, $F(x)=0$, with $m$ distinct roots is called a \emph{cyclic equation} over the field $K$ if its roots $\alpha,\alpha_1,\alpha_2,\dots,\alpha_{m-1}$ are not rational, but can be arranged in such a way that the cyclic functions of the roots are rational over $K$.

If therefore
\[
\pi=(\alpha,\alpha_1,\alpha_2,\dots,\alpha_{m-1}),
\tag{6}
\]
then every function belonging to $K$ and invariant under the permutations of the cyclic group
\[
C=1,\pi,\pi^2,\pi^3,\dots,\pi^{m-1}
\tag{7}
\]
must be rational.

The Galois group of a cyclic equation is either the whole period $C$, in which case the equation is irreducible, or else it is a subgroup of $C$; in the latter case, if $e$ and $f$ are two integer divisors of $m$ with $m=ef$, then this group consists of the permutations
\[
C^e = 1,\pi^e,\pi^{2e},\dots,\pi^{(f-1)e},
\tag{8}
\]
and the cyclic equation $F(x)=0$ splits into $e$ factors of degree $f$.

For let
\[
t=\psi(\alpha,\alpha^e,\alpha^{2e},\dots,\alpha^{(f-1)e})
\]
be a function belonging to the group $C$, and assume it is equal to a quantity in $K$, which we denote by $a$.

Then on the rational equation
\[
F(x)=F(x,a)F(x,a_1)\cdots F(x,a_{e-1})
\]
no permutation not contained in $C$ can act, and consequently the group of the equation can contain no substitutions other than those occurring in $C$.

If now $\pi$ itself occurs in the group of the equation, then the group is identical with $C$, and since $C$ is transitive, the equation is irreducible.

If, however, $\pi^{e'}$ is the smallest power of $\pi$ occurring in the group of the equation, then $C^{e'}$ is the group. If $e'<m$, then $C^{e'}$ is intransitive, and the equation is reducible.

The same reduction can also be applied to cyclic equations whose degree is not prime. If namely $m=ef$ is any factorization of $m$ into two factors, then the permutation $\pi^e$ splits into $e$ cycles $\gamma$, each with $f$ elements:
\[
\gamma=(\alpha,\alpha^e,\alpha^{2e},\dots,\alpha^{(f-1)e}).
\tag{9}
\]
Now determine a function $\psi$ belonging to the first of these cycles, and set
\[
\eta=\eta_0=\psi(\alpha,\alpha^e,\alpha^{2e},\dots,\alpha^{(f-1)e}),
\]
\[
\eta_1=\psi(\alpha_1,\alpha^{e+1},\alpha^{2e+1},\dots,\alpha^{(f-1)e+1}),
\]
\[
\dots,
\]
\[
\eta_{e-1}=\psi(\alpha_{e-1},\alpha_{2e-1},\alpha_{3e-1},\dots,\alpha_{m-1}).
\tag{10}
\]
These quantities are all distinct, yet under application of the permutation $\pi$ they pass cyclically into one another. Hence their cyclic functions, and therefore also their symmetric functions, are rational. So they are the roots of a cyclic equation of degree $e$, while $F(x)$ splits into $e$ factors of degree $f$,
\[
F(x)=F(x,\eta)F(x,\eta_1)\cdots F(x,\eta_{e-1}),
\]
each of which gives a cyclic equation of degree $f$ for the roots of one of the cycles $\gamma$.

Indeed, if for example
\[
F_1=(x-\alpha)(x-\alpha^e)\cdots(x-\alpha^{(f-1)e}),
\]
then this function admits the permutations in the period of $\gamma$. It also admits the permutations of the group that becomes the Galois group of our equation after adjoining $\eta$, a group that can consist only of powers of $\gamma_0,\gamma_1,\dots,\gamma_{e-1}$.

Hence, by Lagrange's theorem (\S 155, equation (2)), $F_1$ is rationally expressible in terms of $\eta$, so $F_1=F(x,\eta)$. But the equation
\[
F(x,\eta)=0
\]
is again cyclic over $K(\eta)$, since the cyclic functions of its roots belong to that field.

Like all Abelian equations, cyclic equations have the property that every root is rationally expressible in terms of every other. Here these expressions can be arranged in the following cyclic manner. If $\alpha,\alpha_1,\alpha_2,\dots,\alpha_{m-1}$ are the roots of the cyclic equation $F(x)=0$, then the entire function of degree $m-1$ in $x$
\[
\Psi_m\!\left(x+\frac{F'(x)}{\Psi'_m(x)}\right)=\Psi_m
\]
is unchanged by the permutation $\pi$ and therefore belongs to $K$. If in it we set
\[
x=\alpha,\alpha_1,\dots,\alpha_{m-1}
\]
and introduce the symbol $\theta$, then it follows that
\[
\alpha_1=\theta(\alpha),\qquad
\alpha_2=\theta(\alpha_1),\qquad
\dots,\qquad
\alpha_{m-1}=\theta(\alpha_{m-2}),\qquad
\alpha=\theta(\alpha_{m-1}),
\tag{11}
\]
and this holds whether $F(x)$ is reducible or irreducible, provided only that $F(x)$ and $F'(x)$ have no common factor.

The solution of cyclic equations of degree $m$ has thus been made dependent on the solution of cyclic equations whose degrees are the prime factors of $m$. If for example $m$ is a power of $2$, then the solution is carried out by a sequence of square roots. In this way Gauss first treated the cyclotomic equations.

Finally we add the remark, important for what follows, that for prime degree the notions of normal equation and cyclic equation coincide. Certainly a cyclic equation of prime degree is irreducible and hence a normal equation.

Conversely, if the degree $n$ of a normal equation, which is also the degree of its group, is prime, and if $\pi$ is any non-identity permutation of this group, then the degree of $\pi$, which must divide $n$, is equal to $n$; hence
\[
\pi^n=1,\qquad \pi,\pi^2,\dots,\pi^{n-1}
\]
constitute the group of the equation, and the equation is cyclic.

\section*{\S 5. Lagrange resolvents}

\subsection*{164. The method of solution for cyclic equations}

The method by which we shall now learn to solve cyclic equations applies uniformly to prime degrees and to composite degrees. One uses for this purpose certain expressions known under the name of \emph{Lagrange resolvents}, which are of great use in all investigations concerning the algebraic solution of equations.

Let $F(x)=0$ be an equation with roots
\[
\alpha,\alpha_1,\alpha_2,\dots,\alpha_{m-1}.
\]
Let $\varepsilon$ be any $m$th root of unity, and introduce the notation
\[
(\varepsilon,\alpha)=\alpha+\varepsilon\alpha_1+\varepsilon^2\alpha_2+\cdots+\varepsilon^{m-1}\alpha_{m-1}.
\tag{1}
\]
These sums are what one calls the \emph{Lagrange resolvents}. If these functions are known for all $m$th roots of unity $\varepsilon$, then the equation itself is solved; for by \S 133, equation (6),
\[
\sum_{\varepsilon}\varepsilon^{-k}=m \quad\text{or}\quad 0,
\tag{2}
\]
according as $k$ is divisible by $m$ or not, and hence
\[
m\alpha_k=\sum_{\varepsilon}(\varepsilon,\alpha)\varepsilon^{-k},
\tag{3}
\]
where the sums extend over all $m$th roots of unity $\varepsilon$.

The other roots can be expressed in the same way:
\[
m\alpha_k=\sum_{\varepsilon}\varepsilon^{-k}(\varepsilon,\alpha),
\tag{4}
\]
so that everything is indeed reduced to knowledge of $(\varepsilon,\alpha)$ and the functions formed from it.

The sums \((2)\) and \((3)\) can also be written in another fashion, since all $m$th roots of unity are powers of a primitive one.

If therefore $\varepsilon$ denotes a fixed primitive $m$th root of unity and $h$ runs through a complete system of residues modulo $m$, say $0,1,\dots,m-1$, then equations \((2)\), \((3)\) may be written as
\[
m\alpha=\sum_h (\varepsilon^h,\alpha),\qquad
m\alpha_k=\sum_h \varepsilon^{-hk}(\varepsilon^h,\alpha).
\tag{5}
\]

We first investigate these resolvents as functions of $m$ independent variables $\alpha$, and for simpler notation we adopt the convention that $\alpha_h=\alpha_k$ whenever $h\equiv k\pmod m$.

Then the entire system of variables $\alpha$ is obtained by letting the index $h$ run through a complete residue system modulo $m$.

For these resolvents the following statements hold.

\medskip
\noindent\textbf{1.}
If one applies to the indices of the $\alpha$ the cyclic permutation
\[
\pi=(0,1,\dots,m-1),
\]
then $(\varepsilon,\alpha)$ passes into $\varepsilon^{-1}(\varepsilon,\alpha)$, and under the permutation $\pi^k$ it passes into $\varepsilon^{-k}(\varepsilon,\alpha)$.

\medskip
This follows immediately from the defining formula \((1)\).

Now let $\nu$ be any positive exponent and form $(\varepsilon,\alpha)^\nu$. Expand this power by the polynomial theorem, set $\varepsilon^m=1$, and order the result by powers of $\varepsilon$. One obtains an expression of the form
\[
(\varepsilon,\alpha)^\nu
=
A^{(\nu)}_0+\varepsilon A^{(\nu)}_1+\varepsilon^2A^{(\nu)}_2+\cdots+\varepsilon^{m-1}A^{(\nu)}_{m-1},
\tag{5'}
\]
where the $A^{(\nu)}_h$ are forms of degree $\nu$ with integer coefficients in the variables $\alpha$, but are independent of $\varepsilon$.
Again we stipulate that $A^{(\nu)}_h=A^{(\nu)}_k$ whenever $h\equiv k\pmod m$.

From this we prove the following statement.

\medskip
\noindent\textbf{2.}
If one applies the cyclic permutation $\pi$ to the indices of the $\alpha$, then the indices of the coefficients in $(\varepsilon,\alpha)^\nu$ undergo the cyclic permutation $\pi^\nu$; that is, $A^{(\nu)}_h$ passes into $A^{(\nu)}_{h+\nu}$.

\medskip
To prove this, note that formula \((5')\) remains valid for every $m$th root of unity $\varepsilon$, including $1$. Using \((2)\) one then obtains
\[
mA^{(\nu)}_k=\sum_{\varepsilon}\varepsilon^{-k}(\varepsilon,\alpha)^\nu.
\tag{6}
\]
If we apply the permutation $\pi$ on the right-hand side, then by statement 1 the term $(\varepsilon,\alpha)^\nu$ is multiplied by $\varepsilon^{-\nu}$, so the whole sum becomes
\[
\sum_{\varepsilon}\varepsilon^{-(k+\nu)}(\varepsilon,\alpha)^\nu
=
mA^{(\nu)}_{k+\nu},
\]
which is exactly the assertion.

\medskip
\noindent\textbf{3.}
Thus, when the permutation $\pi$ is applied to the indices of the $\alpha$, it induces on the indices of the coefficients of the product ordered by powers of $\varepsilon$ the permutation determined by the total exponent-weight.

\medskip
This general principle will be used in the further computation of the Lagrange resolvents.

\medskip

\clearpage

\section*{\S 37. Linear substitutions and their composition}

One of the most effective means for forming groups, and at the same time one that leads to many applications, is given by linear substitutions and their composition. We have already encountered such linear substitutions several times, for example in the second section of the first volume on the occasion of the multiplication law for determinants, and then again in the eleventh section in connection with equivalent numbers.

By a \emph{linear substitution} in \(n\) variables we understand a system of equations by which one system of \(n\) variables \(y_1,y_2,\dots,y_n\) is expressed linearly in terms of another system \(x_1,x_2,\dots,x_n\). According to the number of variables we distinguish unary, binary, ternary, quaternary substitutions, and in general we shall call the number of variables the \emph{dimension} of the substitution. For the time being we restrict attention to homogeneous substitutions, so that, if the coefficients are denoted by \(a^{(k)}_i\), the substitution is represented by the system
\[
y_\kappa=\sum_{i=1}^n a^{(\kappa)}_i x_i
\qquad (\kappa=1,2,\dots,n).
\tag{1}
\]
Very often the variables themselves are not important, so that such a substitution is sufficiently characterized by its coefficients \(a^{(\kappa)}_i\). One sometimes uses a single letter, say \(A\), to denote a substitution, setting
\[
A=
\begin{pmatrix}
a^{(1)}_1 & a^{(1)}_2 & \cdots & a^{(1)}_n \\
a^{(2)}_1 & a^{(2)}_2 & \cdots & a^{(2)}_n \\
\vdots & \vdots & \ddots & \vdots \\
a^{(n)}_1 & a^{(n)}_2 & \cdots & a^{(n)}_n
\end{pmatrix}.
\tag{2}
\]
We shall also write briefly \(A=(a^{(\kappa)}_i)\), and denote the determinant of the substitution coefficients by \(|A|\). This determinant is always assumed to be nonzero. We also write the system symbolically as
\[
(y_1,y_2,\dots,y_n)=A(x_1,x_2,\dots,x_n),
\tag{3}
\]
or, more briefly,
\[
(y)=A(x).
\tag{4}
\]

The substitution
\[
y_1=x_1,\quad y_2=x_2,\quad \dots,\quad y_n=x_n
\tag{5}
\]
or
\[
J=
\begin{pmatrix}
1 & 0 & \cdots & 0\\
0 & 1 & \cdots & 0\\
\vdots & \vdots & \ddots & \vdots\\
0 & 0 & \cdots & 1
\end{pmatrix}
\tag{6}
\]
is called the \emph{identity substitution}.

A substitution of the form
\[
y_1=\mu_1x_1,\quad y_2=\mu_2x_2,\quad \dots,\quad y_n=\mu_nx_n
\tag{7}
\]
or
\[
M=
\begin{pmatrix}
\mu_1 & 0 & \cdots & 0\\
0 & \mu_2 & \cdots & 0\\
\vdots & \vdots & \ddots & \vdots\\
0 & 0 & \cdots & \mu_n
\end{pmatrix}
\tag{8}
\]
will be called a \emph{multiplicative substitution}, or simply a \emph{multiplication}; the elements \(\mu_1,\mu_2,\dots,\mu_n\) are called the \emph{multipliers}.

Now suppose that we have a second linear substitution of the same dimension \(n\), introducing a new system of variables \(z\):
\[
z_i=\sum_{h=1}^n b^{(i)}_h\,y_h,\qquad (z)=B(y).
\tag{9}
\]
If we substitute for \(y\) the values from \((1)\), then we obtain the \(z\) expressed in terms of the \(x\) by means of a new linear substitution \(E\), which we denote by
\[
z=E(x)=BA(x),
\tag{10}
\]
and the substitution \(E=BA\) is said to be \emph{composed} from \(B\) and \(A\). One must, however, distinguish between \(AB\) and \(BA\). Two substitutions \(A,B\) having the special property \(AB=BA\) are called \emph{interchangeable} or \emph{commutative}.

The coefficients of \(E\) are obtained by inserting the expressions \((1)\) into \((9)\):
\[
z_\kappa=\sum_h e^{(\kappa)}_h x_h,\qquad
e^{(\kappa)}_h=\sum_{i=1}^n b^{(\kappa)}_i a^{(i)}_h.
\tag{11}
\]
These formulas are exactly the same as those used earlier for the multiplication of determinants. Thus we may express the composition law of substitutions as follows.

\medskip
\noindent\textbf{1.}
To form the substitution \(BA\) composed from \(B\) and \(A\), one proceeds exactly as if the determinants \(|B|\) and \(|A|\) were to be multiplied by the determinant multiplication rule. In order to obtain the elements of one row, one multiplies the elements of a row of the first factor by the corresponding elements of the columns of the second factor and then adds.

\medskip
From this rule there follows immediately:

\medskip
\noindent\textbf{2.}
The determinant of a composed substitution is equal to the product of the determinants of its components.

\medskip
To compose three substitutions \(C,B,A\) of the same dimension, we insert the expressions \((10)\) for \(z\) into a further linear substitution
\[
(u)=C(z).
\tag{12}
\]
Then the variables \(u\) are expressed in terms of the \(x\) by
\[
(u)=CBA(x).
\tag{13}
\]
Since it is evidently immaterial whether one first introduces \((10)\) into \((12)\), or first expresses \(z\) by \(y\) and then \(y\) by \(x\), one obtains for this composition the associative law
\[
C(BA)=(CB)A=CBA.
\tag{14}
\]
The same may of course be verified directly from \((11)\). Thus:

\medskip
\noindent\textbf{3.}
For the composition of linear substitutions, the associative law holds, but in general not the commutative law.

\medskip
A multiplication in which all multipliers are equal,
\[
N=
\begin{pmatrix}
\nu & 0 & \cdots & 0\\
0 & \nu & \cdots & 0\\
\vdots & \vdots & \ddots & \vdots\\
0 & 0 & \cdots & \nu
\end{pmatrix},
\]
will be called a \emph{similarity substitution}. Two substitutions that are obtained from one another by composition with a similarity substitution, such as \(A\) and \(AN\) or \(A\) and \(NA\), are called \emph{similar}. From the composition law one sees immediately:

\medskip
\noindent\textbf{4.}
A similarity substitution is commutative with every substitution \(A\) of the same dimension; two similarity substitutions, composed with one another, again yield a similarity substitution; and two substitutions that are each similar to a third are also similar to one another.

\medskip
Equally immediate is the statement:

\medskip
\noindent\textbf{5.}
Composition with the identity substitution \(J\) leaves every substitution unchanged.

\medskip
The substitution \(J\) therefore plays the role of a unit in composition and may be omitted wherever it occurs composed with other substitutions.

We have also the following theorem.

\medskip
\noindent\textbf{6.}
To every substitution \(A\) there exists one and only one inverse substitution \(A^{-1}\) satisfying
\[
AA^{-1}=A^{-1}A=J.
\tag{15}
\]

\medskip
This last statement follows from the fundamental formulas of determinant theory if one determines the elements \(\alpha^{(\kappa)}_h\) of \(A^{-1}\) from the linear equations
\[
\sum_i a^{(h)}_i\alpha^{(i)}_\kappa=
\begin{cases}
0,& h\ne \kappa,\\
1,& h=\kappa,
\end{cases}
\tag{16}
\]
from which, if \(A^{(h)}_\kappa\) denotes the minor of \(|A|\), one obtains
\[
|A|\alpha^{(\kappa)}_h=A^{(h)}_\kappa.
\tag{17}
\]
These in turn imply the system
\[
\sum_i \alpha^{(h)}_i a^{(i)}_\kappa=
\begin{cases}
0,& h\ne \kappa,\\
1,& h=\kappa.
\end{cases}
\tag{18}
\]

The inverse substitution of \(A\) is nothing other than the solution of the system of equations defining the direct substitution \(A\), so that from
\[
(y)=A(x)
\]
there follows
\[
(x)=A^{-1}(y).
\tag{19}
\]

For the composition of inverse substitutions one obtains from \(ABB^{-1}A^{-1}=J\) the formula
\[
(AB)^{-1}=B^{-1}A^{-1}.
\tag{20}
\]

If \(A,B,C\) are three substitutions, then composition with \(A^{-1}\) shows from either of the equations
\[
AB=AC,\qquad BA=CA
\]
that necessarily \(B=C\). Hence the characteristic properties required for a group are satisfied by the composition of substitutions, and thus:

\medskip
\noindent
\textbf{The totality of all substitutions of a fixed dimension forms an (infinite) group.}

\clearpage

\section*{Continuation of \S 37. Transformed and transposed substitutions}

If a set of substitutions has the property that the composition of any two of its elements again gives an element of the same set, then that set likewise forms a group, which may be finite or infinite.

Let $L$ be a fixed substitution. From every substitution $A$ of the same dimension one may derive a substitution
\[
A' = L^{-1}AL
\tag{21}
\]
called the \emph{transform} of $A$ by $L$.
If we set
\[
(x)=L(x'),\qquad (y)=L(y'),
\tag{22}
\]
then from the defining relation for $A$ it follows that
\[
(y') = A'(x'),
\tag{23}
\]
so that passing to the transformed substitution is equivalent to applying the same change of variables $L^{-1}$ simultaneously to both systems of variables.

If
\[
A' = L^{-1}AL,\qquad B' = L^{-1}BL,
\]
then the composition laws give
\[
A'B' = L^{-1}ABL,
\]
and therefore:

\medskip
\noindent\textbf{7.}
As $A$ runs through the substitutions of a group, the transformed substitutions $L^{-1}AL$, for fixed $L$, run through the substitutions of an isomorphic group.

\medskip
The inverse substitutions are to be distinguished carefully from the transposed substitutions.

From each substitution $A$ one obtains another specific substitution, called the \emph{transposed substitution}, which for the moment we denote by $A^{t}$, by interchanging rows and columns in $A$ - that is, by swapping the upper and lower indices of the coefficients.

If in the composition formula one interchanges lower and upper indices and at the same time exchanges the roles of $a$ and $b$, one obtains the rule
\[
(BA)^{t} = A^{t}B^{t}.
\tag{24}
\]
Thus the transpose behaves under composition in a way closely analogous to the inverse substitution: in taking transposes one must reverse the order of the factors, whereas for inverses one preserves the order and replaces each factor by its inverse. Combining the two rules gives
\[
(A^{-1})^{t} = (A^{t})^{-1},
\tag{25}
\]
which Weber abbreviates, when no ambiguity is possible, by $A^{-1}_{t}$.

\section*{\S 38. The group of integral linear substitutions of one variable}

We now apply the general results of the preceding paragraph to several special cases and begin with substitutions of one variable. These substitutions are identical with integral fractional-linear functions of the first degree:
\[
y = \frac{\alpha x + \beta}{\gamma x + \delta},
\tag{26}
\]
for which we introduce, for brevity, the symbol
\[
\begin{pmatrix}
\alpha & \beta\\
\gamma & \delta
\end{pmatrix},
\tag{27}
\]
and whose composition is governed by the formula
\[
\begin{pmatrix}
\alpha & \beta\\
\gamma & \delta
\end{pmatrix}
\begin{pmatrix}
\alpha' & \beta'\\
\gamma' & \delta'
\end{pmatrix}
=
\begin{pmatrix}
\alpha\alpha' + \beta\gamma' & \alpha\beta' + \beta\delta'\\
\gamma\alpha' + \delta\gamma' & \gamma\beta' + \delta\delta'
\end{pmatrix}.
\tag{28}
\]
The coefficients $\alpha,\beta,\gamma,\delta$ are arbitrary numbers whose determinant $\alpha\delta-\beta\gamma$ is nonzero.

For the identity substitution one has
\[
\begin{pmatrix}
1 & 0\\
0 & 1
\end{pmatrix},
\tag{29}
\]
and the inverse substitution of \((26)\) is
\[
\frac{1}{D}
\begin{pmatrix}
\delta & -\beta\\
-\gamma & \alpha
\end{pmatrix},
\qquad D=\alpha\delta-\beta\gamma.
\tag{30}
\]

The transposed substitution attached to \((26)\) is obtained by interchanging the diagonal entries and simultaneously changing the signs of the off-diagonal entries:
\[
\begin{pmatrix}
\delta & -\beta\\
-\gamma & \alpha
\end{pmatrix},
\tag{31}
\]
or, what is the same, by replacing $\alpha,\beta,\gamma,\delta$ by $\delta,-\gamma,-\beta,\alpha$. One verifies immediately that
\[
\begin{pmatrix}
\alpha & \beta\\
\gamma & \delta
\end{pmatrix}
\begin{pmatrix}
\delta & -\beta\\
-\gamma & \alpha
\end{pmatrix}
=
\begin{pmatrix}
D & 0\\
0 & D
\end{pmatrix},
\tag{32}
\]
which on the one hand confirms formula \((24)\), and on the other shows that every such substitution commutes with its transpose whenever $D$ is a symmetric function of $\alpha$ and $\delta$.

\section*{\S 39. Permutation groups as linear substitution groups}

Let $x_1,x_2,\dots,x_n$ be a system of $n$ variables, and let $\alpha_1,\alpha_2,\dots,\alpha_n$ be some ordering of the digits $1,2,\dots,n$. Then the equations
\[
x_1=x'_{\alpha_1},\quad x_2=x'_{\alpha_2},\quad \dots,\quad x_n=x'_{\alpha_n}
\tag{33}
\]
determine a linear substitution of dimension $n$,
\[
A=(a_i^{(k)}),\qquad (x)=A(x'),
\tag{34}
\]
in which each row and each column contains exactly one nonzero coefficient, and that coefficient is equal to $1$.

The determinant $|A|$ is therefore equal to $\pm 1$. If, after carrying out the substitution, one sets $x'_i=x_i$ again, the result is simply the permutation
\[
\begin{pmatrix}
1,2,\dots,n\\
\alpha_1,\alpha_2,\dots,\alpha_n
\end{pmatrix}
\]
of the indices of the variables.

If $B$ is a second substitution of the same kind,
\[
(x') = B(x''),
\tag{35}
\]
or explicitly
\[
x'_1=x''_{\beta_1},\quad x'_2=x''_{\beta_2},\quad \dots,\quad x'_n=x''_{\beta_n},
\tag{36}
\]
then the composition rules of \S 37 yield
\[
x_1=x''_{\beta\alpha_1},\quad x_2=x''_{\beta\alpha_2},\quad \dots,\quad x_n=x''_{\beta\alpha_n},
\tag{37}
\]
which is abbreviated by
\[
(x)=AB(x'').
\tag{38}
\]
But according to the composition law for permutations one has
\[
\begin{pmatrix}
1,2,\dots,n\\
\alpha_1,\alpha_2,\dots,\alpha_n
\end{pmatrix}
\begin{pmatrix}
1,2,\dots,n\\
\beta_1,\beta_2,\dots,\beta_n
\end{pmatrix}
=
\begin{pmatrix}
1,2,\dots,n\\
\beta\alpha_1,\beta\alpha_2,\dots,\beta\alpha_n
\end{pmatrix}.
\tag{39}
\]
Thus, if the substitutions $A$ and $B$ are regarded as permutations of the indices, then the substitution $AB$ is exactly the composed permutation $AB$.

Hence permutation groups on $n$ symbols are nothing other than a special case of finite groups of linear substitutions of dimension $n$.

Permutations of the first kind correspond to substitutions with determinant $+1$, and permutations of the second kind correspond to substitutions with determinant $-1$.

This relation between permutations and linear substitutions enables us to derive from every permutation group a group of linear substitutions, and conversely to settle certain questions about permutation groups with the aid of the theory of linear substitutions.

\clearpage

\webersection{40}{Groups of binary linear substitutions with real coefficients}

For the binary substitutions
\begin{equation}
 y_1=\alpha x_1+\beta x_2,\qquad
 y_2=\gamma x_1+\delta x_2,
\end{equation}
the determinant is
\[
D=\alpha\delta-\beta\gamma .
\]
We shall first suppose that the coefficients
\(\alpha,\beta,\gamma,\delta\) are real, and ask under what conditions a group of such substitutions can be finite.

The characteristic equation of the substitution is
\begin{equation}
\begin{vmatrix}
\alpha-\lambda & \beta\\
\gamma & \delta-\lambda
\end{vmatrix}
=
\lambda^2-(\alpha+\delta)\lambda+D=0.
\end{equation}
Its roots are
\begin{equation}
\lambda=
\frac{\alpha+\delta\pm\sqrt{(\alpha+\delta)^2-4D}}{2}.
\end{equation}

If
\[
(\alpha+\delta)^2-4D>0,
\]
then the roots are real and distinct.  In that case the substitution can be brought, by a similar substitution, into the form
\begin{equation}
y_1=\lambda_1x_1,\qquad y_2=\lambda_2x_2.
\end{equation}
The \(n\)-fold iterate is then
\[
y_1=\lambda_1^n x_1,\qquad y_2=\lambda_2^n x_2.
\]
For some positive power to become the identity, \(\lambda_1\) and \(\lambda_2\) must be roots of unity.  Since they are real, this forces
\[
\lambda_1,\lambda_2\in\{1,-1\}.
\]

If
\[
(\alpha+\delta)^2-4D=0,
\]
then the two characteristic roots coincide: \(\lambda_1=\lambda_2=\lambda\).  The substitution can then either be reduced to the diagonal form above, or to the Jordan form
\begin{equation}
y_1=\lambda x_1,\qquad y_2=x_1+\lambda x_2.
\end{equation}
In the latter case the \(n\)-fold iterate contains the term
\[
y_2=n\lambda^{n-1}x_1+\lambda^n x_2,
\]
and therefore cannot be the identity for any positive \(n\), unless the nilpotent part vanishes.  Thus a finite group contains no nontrivial substitution of this Jordan type.

It remains to consider the case
\[
(\alpha+\delta)^2-4D<0.
\]
Then the two characteristic roots are conjugate complex roots.  A real substitution with such roots is, over the real field, similar to a rotation-dilatation.  If the substitution has finite order, its characteristic roots must again be roots of unity; hence the quotient of the two roots is a root of unity and the corresponding projective transformation is periodic.

Consequently, a real binary linear substitution which belongs to a finite group is characterized by the fact that, after multiplication by a scalar if necessary, it can be interpreted as a rotation through an angle commensurable with \(2\pi\), or as one of the exceptional involutory cases in which the characteristic roots are \(1\) and \(-1\).

This gives a first reduction of the problem of finite real binary substitution groups to the problem of finite rotation groups and their associated projective actions.

\webersection{41}{Invariants and covariants of groups of linear substitutions}

The theory of linear substitutions is connected in the closest way with the theory of algebraic forms.  By a \emph{form} in the variables
\[
x_1,x_2,\ldots,x_n
\]
we mean an integral homogeneous function of these variables.  A form
\[
f(x_1,\ldots,x_n)
\]
of degree \(p\) is transformed by a linear substitution \(A\) into another form
\[
f'(x'_1,\ldots,x'_n)
\]
of the same degree in the new variables.  If \(f'=f\), that is, if the form is transformed into itself by the substitution, then \(f\) is called an \emph{invariant} of \(A\).

More generally, suppose that \(f\) is a function not only of the variables \(x\), but also of other variables \(u,v,\ldots\).  If \(f\) is transformed into itself by applying \(A\) to the \(x\)-variables and at the same time applying corresponding substitutions to the variables \(u,v,\ldots\), then \(f\) is called a \emph{covariant}.

We now consider a group \(S\) of linear substitutions in the variables
\[
x_1,\ldots,x_n,
\]
and seek those forms which are simultaneously invariants, or covariants, for all substitutions of the group.

Let
\[
F_1,F_2,\ldots,F_r
\]
be a system of such invariants which are rationally independent, meaning that no identical rational relation holds among them, while every other invariant can be rationally expressed through them.  Such a system is called a \emph{basis} of the invariant system of the group \(S\).

The most general invariant has the form
\begin{equation}
F=A_1F_1+A_2F_2+\cdots+A_rF_r,
\end{equation}
where \(A_1,A_2,\ldots,A_r\) are arbitrary constants, or more generally quantities lying in the appropriate coefficient field.

From the general theory of algebraic forms one obtains the following fundamental facts.

\begin{enumerate}[label=\Roman*.]
\item There always exists a finite system of invariants from which all the remaining invariants may be generated rationally.
\item If the group is finite, the invariant system may be chosen so that it reflects the finite averaging process over the substitutions of the group.
\item The corresponding covariants are obtained in the same manner once the simultaneous transformation law of the auxiliary variables is fixed.
\end{enumerate}

These facts allow the invariant theory of a finite linear substitution group to be treated algebraically: the full system of invariant forms is replaced by a finite basis and the rational relations among its elements.

\clearpage

\webersection{42}{The absolute invariant and the ground field}

Let
\[
F_1,F_2,\ldots,F_r
\]
be a basis of the invariant system \(J\), as determined above, and let \(F\) be any other nonconstant invariant of \(S\).  The forms \(A_1,A_2,\ldots,A_r\) may then be chosen so that
\begin{equation}
F=A_1F_1+A_2F_2+\cdots+A_rF_r.
\end{equation}

Apply all substitutions of the group \(S\) to this identity.  By assumption,
\[
F,F_1,F_2,\ldots,F_r
\]
remain unchanged; the coefficients \(A_1,A_2,A_3,\ldots\), however, may be transformed into
\[
A_1^{(i)},A_2^{(i)},A_3^{(i)},\ldots .
\]
If one adds all identities obtained in this way and divides by the order \(s\) of the group \(S\), then one obtains
\begin{equation}
F=\overline A_1F_1+\overline A_2F_2+\cdots+\overline A_rF_r,
\end{equation}
where now \(\overline A_1,\overline A_2,\ldots\) are themselves invariants of the group \(S\).  By the defining property of the basis, these averaged coefficients are rationally expressible through \(F_1,\ldots,F_r\).

Hence:

\begin{satz}
Every invariant \(F\) can be represented as a rational function of the basis invariants \(F_1,\ldots,F_r\).
\end{satz}

Let
\[
\vartheta_1,\vartheta_2,\ldots,\vartheta_h
\]
be the coefficients of the general substitution of \(S\), and let \(\Omega\) be the rational field generated by these quantities.  Every invariant \(F\) is then a function of the variables \(x\) and of the quantities \(\vartheta\).  If the \(x\)'s are treated as independent variables and the \(\vartheta\)'s as parameters, the totality of all invariants so obtained forms a field.  We shall call it the \emph{field of invariants} of the group \(S\).

The absolute invariants are precisely those functions that remain unchanged under all substitutions of the group and depend only on the orbit of the point under the group action.  Thus they furnish the rational functions on the quotient of the space of variables by the group.

\webersection{43}{The form problem}

The problem of determining the values of
\[
x_1,x_2,\ldots,x_n
\]
which give a prescribed value \(F_0\) to a given invariant \(F\), when the coefficients of the group are assumed known, is called the \emph{form problem} of the group \(S\).

Let \(\Omega(t)\) denote a general function of the field \(\Omega\) which is carried into itself by the substitutions of \(S\).  The roots of the equation
\[
\Omega(t)=\omega,
\]
where \(\omega\) is any fixed value, are all distinct and are transformed into one another by the substitutions of \(S\).

The solution of the form problem is closely connected with the solution of the equation
\[
\Omega(t)=\omega.
\]
Indeed, if \(\vartheta\) is a root of this equation and one sets
\begin{equation}
x_i=\vartheta_i(\vartheta),
\end{equation}
for suitable functions \(\vartheta_i\), then one obtains values of the \(x_i\) satisfying all invariant relations.

If the invariant \(J\) can be represented as a rational function of \(\vartheta\), then this function cannot change when any of the substitutions
\[
(\vartheta,\vartheta_1),\quad
(\vartheta,\vartheta_2),\quad\ldots
\]
is applied.  Consequently it is rationally expressible through the coefficients of \(\Omega(t)\).  Whether this representation is possible depends on the nature of the group.

If \(S'\) is a subgroup of \(S\) of index \(j\), and no larger subgroup has the same invariance property, then \(\vartheta\) satisfies an equation of degree \(j\).  This equation is to be regarded as a \emph{resolvent} of the form problem.  Every function which admits exactly the substitutions of \(S'\) satisfies such a resolvent.

Thus the form problem transforms invariant theory into an equation problem: the passage from a group to its subgroups corresponds to the passage from a general equation to its resolvents.

\clearpage

\webersection{44}{Galois groups and resolvents}

The theory of groups of linear substitutions is closely related to the theory of algebraic equations.  Let
\begin{equation}
t^n+a_1t^{n-1}+a_2t^{n-2}+\cdots+a_n=0
\end{equation}
be an equation of degree \(n\), with roots
\[
\alpha_1,\alpha_2,\ldots,\alpha_n.
\]
The totality of all permutations of these roots which leave unchanged every rational relation among them forms the \emph{Galois group} of the equation.

A function
\[
F(\alpha_1,\ldots,\alpha_n)
\]
of the roots that is unchanged by all permutations of the Galois group is rationally expressible through the coefficients
\[
a_1,\ldots,a_n
\]
of the equation.

Solving the equation is equivalent to determining a function \(\vartheta\) of the roots that is left fixed by no nonidentity permutation of the Galois group.  Such a function is called a \emph{primitive function} of the field.

The Galois group may always be regarded as a group of linear substitutions: replace each permutation of the roots by the corresponding substitution of the variables
\[
x_1,\ldots,x_n,
\]
as described in \S 39.  In this way the group of the equation becomes a concrete linear substitution group, and the methods of invariant theory become available for the study of algebraic equations.

Resolvents arise by passing to subgroups.  If a function of the roots is invariant under a subgroup \(S'\) but not under the whole group \(S\), then the conjugates of that function under \(S\) are indexed by the cosets of \(S'\).  Their elementary symmetric functions are invariant under \(S\), hence rational functions of the coefficients of the original equation.  The function therefore satisfies an auxiliary equation whose degree is the index \([S:S']\).

\webersection{45}{Metacyclic groups}

A group \(S\) is called \emph{metacyclic} if it has a distinguished chain of subgroups
\begin{equation}
S,\; S_1,\; S_2,\;\ldots,\; S_\nu=1,
\end{equation}
such that each \(S_i\) is a normal subgroup of \(S_{i-1}\) and each quotient group
\[
S_{i-1}/S_i
\]
is cyclic.

Such a group can be built by solving a chain of equations, each of which is cyclic.  The solvability of an algebraic equation by radicals is therefore expressed group-theoretically by the metacyclic nature of its Galois group, or equivalently by the existence of a suitable chain with cyclic factors.

For the symmetric group on \(n\) letters, metacyclicity occurs only for \(n\leq 4\).  Likewise the alternating group on \(n\) letters is metacyclic only for \(n\leq 4\).  This is the group-theoretic obstruction behind the impossibility, in general, of solving equations of degree at least five by radicals.

\webersection{46}{Abstract and concrete groups}

Until now the exposition of group theory has been tied to definite operations: substitutions, linear transformations, and similar concrete processes.  But the group concept may also be detached from any particular representation and defined purely formally.

By an \emph{abstract group} one understands a system of elements
\[
A,B,C,\ldots
\]
for which a composition law, or multiplication, is defined and satisfies the following rules.

\begin{enumerate}[label=\arabic*.]
\item From any two elements \(A,B\), their product is a uniquely determined element \(C=AB\).
\item The associative law holds:
\[
(AB)C=A(BC).
\]
\item There is an \emph{identity element} \(J\) such that
\[
JA=AJ=A
\]
for every element \(A\).
\item For every element \(A\) there is an inverse element \(A^{-1}\) such that
\[
AA^{-1}=A^{-1}A=J.
\]
\end{enumerate}

Two abstract groups are called \emph{isomorphic} if their elements can be named in such a way that the multiplication rule is the same in both groups.  This definition separates the intrinsic structure of the group from the particular objects on which the group happens to act.

The point of abstraction is not to discard concrete examples, but to recognize when two different systems of operations have the same internal law of composition.  Permutation groups, linear substitution groups, rotation groups, and Galois groups can then be compared by their multiplication tables rather than by their external form.

\clearpage

\webersection{47}{The alternating group and its invariants}

The alternating group on \(n\) letters consists of all even permutations and has index \(2\) in the symmetric group.  The function
\begin{equation}
\Delta=\prod_{i<k}(x_i-x_k)
\end{equation}
is an invariant of the alternating group.  More precisely, it is an invariant which changes sign under every odd permutation.

The square
\[
\Delta^2
\]
is therefore an absolute invariant of the symmetric group.  It is a symmetric function of the \(x_i\), and is represented by the discriminant of the corresponding equation.

The most general invariant of the alternating group can be written in the form
\begin{equation}
F=F_1+\Delta F_2,
\end{equation}
where \(F_1\) and \(F_2\) are symmetric functions, that is, invariants of the full symmetric group.

For \(n\geq 5\), the alternating group is simple: it has no normal subgroup other than the identity and itself.  This fact is one of the central structural reasons why the general equation of degree at least five cannot be solved by radicals.

\webersection{48}{Orthogonal substitutions}

A linear substitution is called \emph{orthogonal} if it transforms the quadratic form
\begin{equation}
x_1^2+x_2^2+\cdots+x_n^2
\end{equation}
into itself.  For a substitution \(A\) with coefficients \(a_i^{(k)}\), this means that the identity
\begin{equation}
\sum_k\left(\sum_i a_i^{(k)}x_i\right)^2=\sum_i x_i^2
\end{equation}
holds in the variables \(x\).

Comparing coefficients gives the relations
\begin{equation}
\sum_k a_i^{(k)}a_j^{(k)}
=
\begin{cases}
1,& i=j,\\
0,& i\neq j.
\end{cases}
\end{equation}

For example, if
\[
A=
\begin{pmatrix}
a_1&b_1&c_1\\
a_2&b_2&c_2\\
a_3&b_3&c_3
\end{pmatrix},
\]
then the orthogonality relations are
\begin{align}
a_1^2+b_1^2+c_1^2&=1,&
a_1a_2+b_1b_2+c_1c_2&=0,\\
a_2^2+b_2^2+c_2^2&=1,&
a_1a_3+b_1b_3+c_1c_3&=0,\\
a_3^2+b_3^2+c_3^2&=1,&
a_2a_3+b_2b_3+c_2c_3&=0.
\end{align}
These equations say exactly that the row vectors of \(A\) form an orthonormal system.

From these relations one obtains
\[
|A|^2=1,
\]
and hence
\[
|A|=\pm1.
\]
The orthogonal substitutions with determinant \(+1\) are those of the \emph{first kind}, or rotations.  Those with determinant \(-1\) are those of the \emph{second kind}, or reflections.

The totality of all orthogonal substitutions in \(n\) variables forms a group, the \emph{orthogonal group}.  One may regard either the positions of the body or the rotations themselves as the elements of the group.

\webersection{49}{Linear fractional substitutions}

We shall now show that the group of ternary orthogonal substitutions is isomorphic to the group of linear fractional substitutions in one variable, or equivalently to the group of binary substitutions acting on ratios.

Write a ternary orthogonal substitution as
\begin{equation}
(y_1,y_2,y_3)=A(x_1,x_2,x_3),
\end{equation}
so that
\begin{equation}
y_1^2+y_2^2+y_3^2=x_1^2+x_2^2+x_3^2.
\end{equation}

We first apply a fixed, nonorthogonal change of variables of the form
\begin{align}
x_1&=\xi_1+\xi_2,&
y_1&=\eta_1+\eta_2,\\
x_2&=\xi_1-\xi_2,&
y_2&=\eta_1-\eta_2,\\
x_3&=\xi_3,&
y_3&=\eta_3.
\end{align}
Under this transformation the quadratic condition is converted into an equivalent condition in the new variables.

Now introduce binary variables by setting, in effect,
\[
x_1=\sqrt2\,\xi_1\xi_3,\qquad
x_2=\sqrt2\,\xi_2\xi_3.
\]
Then the corresponding quadratic relation vanishes identically.  The transformed variables \(y_1,y_2,y_3\) become binary quadratic forms in \(\xi_1,\xi_2\), and these forms must satisfy the transformed quadratic equation identically.

The product relation then forces the relevant quadratic forms to decompose into linear factors.  Since \(y_2\) and \(y_3\) cannot have a common factor without making the determinant of the transformation vanish, the factors are arranged in such a way that the ternary orthogonal substitution is represented by a binary linear substitution, considered projectively.

Thus ternary rotations are represented by fractional linear transformations.  In modern language, this is the classical passage from the binary linear group, modulo scalar multiplication, to the rotation group in three dimensions.

\webersection{51}{Finite groups of linear fractional substitutions; poles of groups}

From the preceding developments it follows that, once all finite groups of linear fractional substitutions are known, one immediately obtains all finite groups of proper ternary orthogonal substitutions and all finite groups of binary linear substitutions with determinant \(1\).

Let \(S\) be a finite group of linear fractional substitutions in a variable \(x\).  Write its substitutions in canonical form
\begin{equation}
x'=\frac{\alpha x+\beta}{\gamma x+\delta},
\end{equation}
with determinant condition
\[
\alpha\delta-\beta\gamma=1.
\]

A nonidentity substitution of this form generally has two fixed points, or \emph{poles}, obtained from the quadratic equation
\begin{equation}
\gamma x^2+(\delta-\alpha)x-\beta=0.
\end{equation}
Only in the exceptional case \(\gamma=0\) and \(\alpha=\delta\), that is, for a translation \(x'=x+\beta\), does the equation reduce to degree one; then there is only one improper pole, situated at infinity.

The pole structure is the key to the classification of finite fractional linear groups.  A finite group acts on the sphere of projective values, and the exceptional points with nontrivial stabilizer determine the group.

\webersection{52}{The cyclic and dihedral groups}

The simplest finite groups are the \emph{cyclic groups}.  A cyclic group of degree \(n\) is generated by the substitution
\begin{equation}
x'=\varepsilon x,
\end{equation}
where \(\varepsilon\) is a primitive \(n\)-th root of unity.  The \(n\) substitutions of the group are
\begin{equation}
x,\;\varepsilon x,\;\varepsilon^2x,\;\ldots,\;\varepsilon^{n-1}x.
\end{equation}

The two poles are \(0\) and \(\infty\), and both are fixed by every substitution of the group.

A \emph{dihedral group} of degree \(2n\) contains, besides these substitutions, the \(n\) substitutions
\begin{equation}
x'=\frac{\varepsilon^k}{x},
\qquad k=0,1,\ldots,n-1.
\end{equation}
These interchange the two poles \(0\) and \(\infty\).  Thus the group consists of the substitutions
\begin{equation}
x'=\varepsilon^k x,\qquad
x'=\frac{\varepsilon^k}{x},
\qquad k=0,1,\ldots,n-1.
\end{equation}

The element \(\varepsilon\) must run through a full system of \(n\)-th roots of unity.  Hence, if the \(\varepsilon\) occurring in the generator is primitive, the whole cyclic subgroup is obtained from its powers, and the dihedral group is obtained by adjoining the inversion \(x\mapsto 1/x\), together with its products with the rotations.

\end{document}
